\section{Local bounds from the planted law}\label{sec:planted}

We now bound the center integral for fixed point Lagrangians, including
triples whose pair intersections are nonzero.  Throughout this section the
ambient symplectic space $\mathcal V$ has dimension $2D$, where $D=2r$.
We use the pair probability $\pi_{N,h}$ and the unsigned normalized
incidence kernel $J$ from Section~\ref{sec:orbits}.  Thus
\[
 J(L,Y)=\frac{\1_{\{\dim(L\cap Y)=r\}}}{\pi_{D,r}},
 \qquad \pi_{N,h}\asymp q^{-b(h)}.
\]
For a set $A\subseteq\{1,2,3\}$ of active incidences, put
\begin{equation}\label{eq:planted-face-factor}
 \mathcal F_A(L_1,L_2,L_3)
   =\E_Y\prod_{i\in A}J(L_i,Y),
\end{equation}
where $Y$ is a uniform Lagrangian.  Constants in this section may depend
on the fixed dimension $D$, but are uniform over the subspaces, their
orbits, and the odd primes $q$.  Every exponent displayed below is an
explicit expression in the dimensions; there is no unspecified error
term in an exponent.

\subsection{The density ratio and the common intersection}

Let $\mu$ be the law of three independent uniform Lagrangians.  Define
the \emph{planted law} $\nu$ by first choosing $Y$ uniformly and then,
conditional on $Y$, choosing the three $L_i$ independently and uniformly
subject to $\dim(L_i\cap Y)=r$.  Pair transitivity makes the conditioning
probability equal to $\pi_{D,r}$ for every $Y$.  Hence, exactly,
\begin{equation}\label{eq:planted-density}
 \frac{d\nu}{d\mu}(L_1,L_2,L_3)
 =\pi_{D,r}^{-3}\E_Y
       \prod_{i=1}^3\1_{\{\dim(L_i\cap Y)=r\}}
 =\mathcal F_{\{1,2,3\}}(L_1,L_2,L_3).
\end{equation}
The right side is invariant under simultaneous symplectic maps.
Therefore, on any actual ordered symplectic orbit $\mathcal O$,
\begin{equation}\label{eq:planted-orbit-ratio}
 \mathcal F_{\{1,2,3\}}\big|_{\mathcal O}
     =\frac{\nu(\mathcal O)}{\mu(\mathcal O)}.
\end{equation}
It is essential here to use the mass of this orbit in the denominator.
If its common intersection has dimension $c$, its opposite pair
intersection dimensions are $h_i$, and
\begin{equation}\label{eq:planted-local-dimensions}
 s_i=h_i-c,\qquad d=D-c,\qquad
 C=(D+1)c+\sum_{i=1}^3 b(s_i),
\end{equation}
then Proposition~\ref{prop:triple-orbits} gives
$\mu(\mathcal O)\asymp q^{-C}$, uniformly also when the residual
transverse space has dimension zero.

We estimate the numerator through four geometric choices. First record
the actual common space \(C_0=L_1\cap L_2\cap L_3\), then count how
the center \(Y\) meets it. In the quotient by \(C_0\), count the
intersections of the three point Lagrangians with the reduced center.
Finally count the remaining lift forms, which supply any additional
pair intersections. These counts use the independent laws before the
mutual point profile is imposed.

The common space \(C_0\) is unique. Thus summing over
isotropic $c$-spaces $C_0$, while retaining that equality, counts each
triple of common dimension $c$ exactly once.  For a deterministic
summand $C_0$, we evaluate probabilities under the original law of
four independent uniform Lagrangians.  Under this law the
three containments $C_0\subseteq L_i$ have total probability
$\asymp q^{-3Dc+3\binom c2}$.  There are
$\asymp q^{2Dc-(3c^2-c)/2}$ choices of $C_0$, so their combined
contribution is
\begin{equation}\label{eq:planted-common-count}
 q^{2Dc-(3c^2-c)/2-3Dc+3\binom c2}
      =q^{-(D+1)c}
\end{equation}
up to a constant factor.  These are the counts of
Lemma~\ref{lem:lagrangian-counts}.

For the center let $u=\dim(C_0\cap Y)$.  Choosing a prospective
$u$-space in $C_0$ and requiring its containment in $Y$ gives
\begin{align}
 \P\{\dim(C_0\cap Y)=u\}
 &\le O\bigl(q^{u(c-u)-Du+\binom u2}\bigr) \notag\\
 &=O\bigl(q^{-u(D-c)-b(u)}\bigr).
 \label{eq:planted-center-common-cost}
\end{align}
This bound includes $u=0$.  We now work in the symplectic quotient
$C_0^\perp/C_0$ and write
\[
 L_i'=L_i/C_0,\qquad
 Y'=\bigl((Y\cap C_0^\perp)+C_0\bigr)/C_0.
\]
Both are Lagrangians of dimension $d$.  There is an exact quotient
identity
\begin{equation}\label{eq:planted-quotient-intersection}
 L_i'\cap Y'=\bigl((L_i\cap Y)+C_0\bigr)/C_0,
 \qquad \dim(L_i'\cap Y')=\dim(L_i\cap Y)-u.
\end{equation}
Indeed, a class in the left side has a representative
$y\in Y\cap C_0^\perp$.  Membership in $L_i/C_0$ means that
$y+z\in L_i$ for some $z\in C_0$.  Since $C_0\subseteq L_i$,
this implies $y\in L_i\cap Y$.  The converse is immediate, and the
map from $L_i\cap Y$ to the quotient has kernel $C_0\cap Y$.
On the required incidence event the last dimension is $r-u$.
This also proves that no choice of lifts is missing from the quotient
count.

Put
\begin{equation}\label{eq:planted-reduced-parameters}
 r'=r-u,\qquad \delta=d-2r'=2u-c.
\end{equation}
For this counting step fix deterministic $C_0,Y$ with
$\dim(C_0\cap Y)=u$.  Under the original independent product law,
condition only on the three containments $C_0\subseteq L_i$.
The quotient Lagrangians $L_i'$ are then independent and uniform,
because containment reduction is a bijection.  The equality
$C_0=\bigcap_i L_i$ is still an event to count.  Their three
individual conditions $\dim(L_i'\cap Y')=r'$ consequently have
probability $\pi_{d,r'}^3\asymp q^{-3b(r')}$.  Including the planted
normalizer $\pi_{D,r}^{-3}$, the gain relative to
Equation~\eqref{eq:planted-common-count} is
\begin{align}
 g_C
 &=3\bigl[b(r)-b(r-u)\bigr]-u(D-c)-b(u) \notag\\
 &=\frac32(2ru-u^2+u)-u(2r-c)-\frac12(u^2+u)
      \notag\\
 &=(r+c+1)u-2u^2.
 \label{eq:planted-common-gain}
\end{align}

The preceding independent law is the law \emph{before} imposing the
mutual intersection profile of the $L_i'$.  In particular, their
common-zero condition is still an event to be counted.  Under their
three individual center conditions the spaces
\[
 R_i=L_i'\cap Y'
\]
are independent uniform $r'$-subspaces of $Y'$.  Conditional on these
individual $R_i$, the remaining choices of $L_i'$ are independent as
well.  We use these product laws below, and never assert independence
after conditioning on the mutual profile.

\subsection{The subspace cost and the lift forms}

The common-space calculation accounts for how the center meets $C_0$.
It remains to count the prescribed pair intersections in the quotient:
first inside $Y'$, and then through the remaining choices of the $L_i'$.

Because the required quotient triple has common intersection zero,
the $R_i$ must have common intersection zero.  Let
\begin{equation}\label{eq:planted-subspace-parameters}
 \begin{split}
 p_i&=\dim(R_j\cap R_k),\qquad t_i=s_i-p_i,\qquad
 P=\sum_i p_i,\qquad T=\sum_i t_i,\\
 v&=d-\dim(R_1+R_2+R_3),\qquad
 w=3r'-d-P+v=\frac{d-3\delta}{2}-P+v,
 \end{split}
\end{equation}
where $\{i,j,k\}=\{1,2,3\}$.
The pair spaces $R_j\cap R_k$ form a direct sum.  To see this, in a
relation $a_1+a_2+a_3=0$ among the three pair spaces, membership of
two terms in $R_i$ forces the remaining term into $R_i$ as well;
each term therefore lies in the common intersection and vanishes.
The related inequality $w\ge0$ follows in the abstract direct sum
$R_1\oplus R_2\oplus R_3$.  The summation map has kernel of
dimension $3r'-d+v$.  The three pair-relation spaces consist of
vectors of the forms
\[
 (0,a,-a),\qquad (-b,0,b),\qquad (z,-z,0).
\]
Their sum can vanish only if $a=b=z$, hence only if all three vanish.
They therefore contribute $P$ independent dimensions to that kernel.

Here $p_i$ records the pair-intersection dimension already inside $Y'$,
and $t_i=s_i-p_i$ is the target additional dimension from the remaining
Lagrangian choices. The parameter $v$ is the span codimension in $Y'$;
$w$ is the dimension of the summation-map kernel after quotienting by
the subspace generated by the pair relations.

\begin{lemma}\label{lem:planted-subspace-cost}
For three independent uniform $r'$-subspaces of a $d$-space, the
event of common intersection zero, opposite pair dimensions $p_i$,
and span codimension $v$ has probability
\begin{equation}\label{eq:planted-subspace-cost}
 O\left(q^{-\delta P-\sum_i p_i^2-vw}\right).
\end{equation}
\end{lemma}
\begin{proof}
Write $n=d-v$.  First choose a prospective span $Z$ of dimension
$n$.  There are $\asymp q^{nv}$ choices, and the probability that
all three independent $r'$-subspaces lie in it is
$\asymp q^{-3r'v}$.  Conditional on these separate containments,
the three subspaces are independent uniform $r'$-subspaces of $Z$.

The number of choices of three pair spaces in direct sum inside
$Z$ is $O(q^{nP-\sum_i p_i^2})$.  Indeed, choose their $P$-dimensional
sum and then an ordered direct-sum decomposition; the two exponents
are $P(n-P)$ and $P^2-\sum_i p_i^2$.  Each $R_i$ must contain the
sum of the other two pair spaces, of dimension $P-p_i$.  In $Z$
this costs exponent $(n-r')(P-p_i)$, for a total of
$2(n-r')P$.  We may drop the requirements that $Z$ be exactly the
span and that these containments give exactly the pair intersections.
Thus the total cost exponent is
\begin{align*}
 (3r'-d+v)v-nP+\sum_i p_i^2+2(n-r')P
 &= (3r'-d+v)v+(d-v-2r')P+\sum_i p_i^2\\
 &=\delta P+\sum_i p_i^2+v(3r'-d-P+v),
\end{align*}
which proves the claim.  Infeasible containments have probability zero.
\end{proof}

We next describe the remaining Lagrangian choices without any
generic-position assumption.  Choose symplectic coordinates
$Y'\oplus(Y')^*$ with pairing
\[
 \omega((y,\alpha),(z,\beta))=\beta(y)-\alpha(z).
\]
For an $r'$-subspace $R\subseteq Y'$, let $A=R^\circ\subseteq(Y')^*$
be its annihilator.  A Lagrangian $L'$ meeting $Y'$ exactly in $R$
projects into $A$ by orthogonality to $R$; its kernel is $R$,
so the dimensions show that its image is all of $A$.  Its induced map
$A\longrightarrow Y'/R\cong A^*$ is a symmetric bilinear form
$F$ on $A$.  Conversely, every such form gives exactly one lift,
namely
\begin{equation}\label{eq:planted-lift-parametrization}
 L'(R,F)=\{(y,\alpha):\alpha\in A,\quad
              \beta(y)=F(\alpha,\beta)\text{ for every }\beta\in A\}.
\end{equation}
The set in Equation~\eqref{eq:planted-lift-parametrization} has dimension
$d$, intersects $Y'$ in $R$, and is isotropic precisely because $F$
is symmetric.  This proves the bijection, including when $A=0$.
There are exactly $q^{b(d-r')}$ lifts for each $R$.  Under the
product law already specified, conditional on the $R_i$, the lift
forms $F_i$ on $A_i=R_i^\circ$ are therefore independent uniform
symmetric forms.

For the opposite pair put
\begin{equation}\label{eq:planted-pair-annihilators}
 Q_i=A_j\cap A_k=(R_j+R_k)^\circ,
 \qquad \dim Q_i=d-2r'+p_i=\delta+p_i.
\end{equation}
Use cyclic differences
\[
 D_1=(F_2-F_3)|_{Q_1},\qquad
 D_2=(F_3-F_1)|_{Q_2},\qquad
 D_3=(F_1-F_2)|_{Q_3}.
\]
Projection of $L_j'\cap L_k'$ to $(Y')^*$ has kernel
$R_j\cap R_k$ and image $\operatorname{rad}D_i$.  In fact, a vector
$\alpha\in Q_i$ has simultaneous lifts exactly when the two
prescribed classes in $Y'/R_j$ and $Y'/R_k$ agree in
$Y'/(R_j+R_k)$.  Pairing this agreement against every
$\beta\in Q_i$ is exactly the condition
$(F_j-F_k)(\alpha,\beta)=0$.  Thus there is an exact sequence
\begin{equation}\label{eq:planted-pair-exact-sequence}
 \begin{gathered}
 0\longrightarrow R_j\cap R_k
 \longrightarrow L_j'\cap L_k'
 \longrightarrow\operatorname{rad}D_i\longrightarrow0,\\
 \dim(L_j'\cap L_k')=p_i+\operatorname{nullity}D_i.
 \end{gathered}
\end{equation}
Thus the required reduced pair profile is exactly the event
$\operatorname{nullity}D_i=t_i$ for all $i$.

The remaining task is therefore to bound the simultaneous nullities of
the three difference forms. We first determine their exact joint law,
rather than assuming independence.

\begin{lemma}\label{lem:planted-difference-image}
Let $A_1,A_2,A_3$ be arbitrary subspaces of a vector space, put
$Q_i=A_j\cap A_k$, and put $V=A_1\cap A_2\cap A_3$.
The image of the cyclic difference map on symmetric forms is exactly
\begin{equation}\label{eq:planted-difference-image}
 \left\{(D_1,D_2,D_3):
       D_1|_V+D_2|_V+D_3|_V=0\right\}.
\end{equation}
Its codimension in the product of the three form spaces on the
$Q_i$ is $b(\dim V)$.  Independent uniform $F_i$ give the uniform
law on this image.
\end{lemma}
\begin{proof}
The displayed constraint is necessary by cancellation.  Conversely,
take any three target forms satisfying it, and prescribe the values
of the three original forms on $V$ to be
\[
 H_1=0,\qquad H_2=-D_3|_V,\qquad H_3=D_2|_V.
\]
They have the required three cyclic differences on $V$.  On each
$Q_i$ choose any symmetric extension of the prescribed $H_j$ for
the positive term of its cyclic difference, and prescribe the
negative term there by subtracting $D_i$.

Each $F_i$ has now received prescriptions on exactly two subspaces
of $A_i$.  Their intersection is $V$, and the prescriptions agree
there.  Two symmetric forms on subspaces $B,C$ agreeing on
$B\cap C$ always extend simultaneously: write
$B+C=(B\cap C)\oplus B_0\oplus C_0$, retain the prescribed blocks,
choose the mixed $B_0$--$C_0$ block arbitrarily, and then extend to
the ambient space.  Apply this construction separately for each
$F_i$.  It proves sufficiency even if the three $A_i$ or $Q_i$
have further linear relations.

Finally, the map taking a target triple to the sum of its restrictions
on $V$ is onto the form space on $V$: extend any prescribed form
from $V$ to one $Q_i$ and take the other two forms to be zero.
Its kernel therefore has codimension $b(\dim V)$.  A linear map
between finite vector spaces sends uniform measure to uniform measure
on its image, proving the last assertion.
\end{proof}

In our application $V=(R_1+R_2+R_3)^\circ$ has dimension $v$.
Consequently the difference law is exactly the law of three independent
uniform target forms conditioned on the constraint in
Equation~\eqref{eq:planted-difference-image}, an event of probability
$q^{-b(v)}$.  A uniform symmetric form on an $m$-space has nullity
at least $t$ with probability $O(q^{-b(t)})$.  For completeness,
there are $O(q^{t(m-t)})$ prospective radical $t$-spaces, and
vanishing on one imposes $mt-\binom t2$ independent coefficients;
subtracting gives $-b(t)$.  This also covers $t=m=0$.
Under the independent target law we may multiply the three bounds;
conditioning costs at most $q^{b(v)}$.  We obtain, uniformly for
each actual tuple $(R_1,R_2,R_3)$,
\begin{equation}\label{eq:planted-nullity-cost}
 \P\{\operatorname{nullity}D_i=t_i\ (1\le i\le3)\mid R_1,R_2,R_3\}
   \le O\left(q^{-\sum_i b(t_i)+b(v)}\right).
\end{equation}
An upper bound larger than one is harmless.  The common-zero
condition on the full lifts, or their membership in a particular
ordered orbit, can only decrease this probability.  In particular,
we do not replace those extra conditions by an independence claim.

\subsection{The resulting gain}

For each fixed parameter tuple, the preceding counts bound its
contribution to the numerator in Equation~\eqref{eq:planted-orbit-ratio}
by the product of the following factors, up to a constant depending
on the fixed dimension:
\begin{center}
\begin{tabular}{@{}ll@{}}
\toprule
Choice & Upper-bound factor\\
\midrule
Common space and its three containments
  & \(q^{-(D+1)c}\)\\
Center intersection and normalized incidence conditions
  & \(q^{g_C}\)\\
Subspaces \(R_i\) inside the reduced center
  & \(q^{-\delta P-\sum_i p_i^2-vw}\)\\
Nullities of the restricted difference forms
  & \(q^{-\sum_i b(t_i)+b(v)}\)\\
\bottomrule
\end{tabular}
\end{center}
The last factor includes the possible loss \(q^{b(v)}\) from
conditioning on the shared restriction constraint. We sum over the
finitely many tuples and divide by an individual orbit mass bounded
below by a constant times \(q^{-(D+1)c-\sum_i b(s_i)}\). This cancels
the common-space cost and gives the following pointwise gain.

\begin{proposition}[Three active incidences]\label{prop:planted-gain}
For a local profile $(h_1,h_2,h_3,c)$, we choose one pointwise gain
exponent independent of \(q\) as follows. Let \(\mathcal U\) be the
union, over all odd prime fields and all ordered orbits of this profile,
of the parameter tuples arising from actual quadruples
$(L_1,L_2,L_3,Y)$ with the stated point profile and
$\dim(L_i\cap Y)=r$, using
Equations~\eqref{eq:planted-reduced-parameters} and
\eqref{eq:planted-subspace-parameters}. Their integer entries are bounded
at fixed \(D\), so \(\mathcal U\) is finite. If \(\mathcal U\) is empty,
the face factor is zero over every odd prime field. Otherwise the face
factor with three active incidences is
$O(q^g)$ uniformly over all triples in the profile, where
\begin{equation}\label{eq:planted-gain}
 g=\max_{\text{tuple}\in\mathcal U}\left\{
 g_C+\sum_i b(s_i)-\delta P-\sum_i p_i^2-vw
                   -\sum_i b(t_i)+b(v)\right\},
\end{equation}
so the selected exponent is independent of \(q\).

Every tuple in \(\mathcal U\) satisfies the following necessary
constraints:
\begin{equation}\label{eq:planted-feasibility}
 \begin{gathered}
 0\le u\le\min(c,r),\quad 0\le r'=r-u\le d,
 \quad |\delta|\le c,\quad d-\delta=2r',\\
 P+T=\sum_i s_i\le d,\qquad
 0\le t_i\le p_i+\delta,\qquad
 0\le v\le p_i+\delta\quad(1\le i\le3),\qquad w\ge0.
 \end{gathered}
\end{equation}
In particular $T\le P+3\delta$ and $d+3\delta\ge0$.
These conditions are not claimed sufficient for realizability.
\end{proposition}
\begin{proof}
Fix an actual ordered orbit $\mathcal O$.  Start with four independent
ambient Lagrangians and the normalizer $\pi_{D,r}^{-3}$ in
Equation~\eqref{eq:planted-density}.  Partition its contribution
according to the unique common space $C_0$, and then according to
$u,p_i,t_i,v$.  Equations~\eqref{eq:planted-common-count}--\eqref{eq:planted-common-gain},
Lemma~\ref{lem:planted-subspace-cost}, and
Equation~\eqref{eq:planted-nullity-cost} give the numerator bound
\[
 \nu(\mathcal O)\le
 O\left(\sum_{\text{actual parameter tuples}}
 q^{-(D+1)c+g_C-\delta P-\sum_i p_i^2-vw
                      -\sum_i b(t_i)+b(v)}\right).
\]
At intermediate counting steps we have discarded exactness or orbit
conditions only to enlarge the count.  The parameter tuples indexing
the original partition remain actual ones.  Dividing by the lower
bound $\mu(\mathcal O)\gg q^{-(D+1)c-\sum_i b(s_i)}$ from
Proposition~\ref{prop:triple-orbits} proves
Equation~\eqref{eq:planted-gain}.  There are only boundedly many
dimension tuples for fixed \(D\). The tuples from the current field
belong to \(\mathcal U\), so its maximum gives a
bound uniform in \(q\).

The bounds on $u,r',\delta$ follow from their definitions and
Equation~\eqref{eq:planted-quotient-intersection}.  The reduced pair
intersections are in isotropic direct sum, giving
$\sum_i s_i\le d$ as in Proposition~\ref{prop:triple-orbits}.
Equation~\eqref{eq:planted-pair-exact-sequence} gives the bounds on
$t_i$, and $V\subseteq Q_i$ gives those on $v$.  We already proved
$w\ge0$.  Summing $t_i\le p_i+\delta$ gives
$T\le P+3\delta$, and summing $v\le p_i+\delta$ gives
$v\le P/3+\delta$.  Thus $P+3\delta\ge0$ and
\[
 d+3\delta\ge 6\delta+2P-2v
       \ge\frac43(P+3\delta)\ge0,
\]
where the first inequality uses $w\ge0$.  There are further
necessary conditions, such as $p_i\le r'$ and $P-p_i\le r'$.
The later estimates need only Equation~\eqref{eq:planted-feasibility};
maximizing over actual tuples preserves all the others automatically.
\end{proof}

The gain bound separates three sources of intersection: $u$ records
how the center meets the common space, $p_i$ the pair intersections
inside the reduced center, and $t_i$ the additional intersections
supplied by lift forms. Their costs give one pointwise exponent for
the center factor. The global comparison also needs the cases with
fewer active incidences.

\begin{lemma}[At most two active incidences]\label{lem:two-active-gain}
For two active incidences whose point Lagrangians meet in dimension
$h$, a valid gain exponent is
\begin{equation}\label{eq:planted-two-active-gain}
 g_2(h)
 =\max_{\substack{p+t=h\,;\ 0\le t\le p\le r}}
      \{b(h)-p^2-b(t)\}
 =b(\lfloor h/2\rfloor).
\end{equation}
For zero or one active incidence, the face factor is exactly one.
\end{lemma}
\begin{proof}
Use the two-point version of the planted density identity.  Under
this law the two $R_i=L_i\cap Y$ are independent uniform
$r$-subspaces of the $2r$-space $Y$.  Their intersection dimension
$p$ has probability $O(q^{-p^2})$: for fixed $R_1$, choose a
$p$-subspace of it, in $O(q^{p(r-p)})$ ways, and require its
containment in $R_2$, at cost $q^{-pr}$.  Conditional on the two
individual $R_i$, the two lift forms are independent uniform.  Their
difference on $(R_1+R_2)^\circ$, of dimension $p$, is a uniform
symmetric form, since restriction from either original form space
is surjective.  Equation~\eqref{eq:planted-pair-exact-sequence}
now gives $h=p+t$ with $0\le t\le p$, and the nullity event costs
at most $O(q^{-b(t)})$.  The ambient pair orbit has mass
$\asymp q^{-b(h)}$ by Lemma~\ref{lem:lagrangian-counts}.  Dividing
gives the maximand in Equation~\eqref{eq:planted-two-active-gain}.

For fixed $h$, let $f(p)=b(h)-p^2-b(h-p)$.  On its integer domain
$\lceil h/2\rceil\le p\le\min(h,r)$,
\[
 f(p+1)-f(p)=h-3p-1<0.
\]
Its maximum is therefore at $p=\lceil h/2\rceil$, which is
allowed since $h\le2r$.  If $h=2a$ the value is
$b(2a)-a^2-b(a)=b(a)$; if $h=2a+1$ it is
$b(2a+1)-(a+1)^2-b(a)=b(a)$.  The maximizing parameters are
realizable: take two $r$-subspaces with intersection $p$, then a
symmetric difference form of nullity $t=h-p$.  Pair transitivity
makes this applicable to every pair of intersection dimension $h$.
Finally, an empty product is one, and for one incidence the mean-one
normalization of $J(L,Y)$ gives exactly one for every $L$.
\end{proof}
